\subsubsection{Transformer Fine-Tuning}

Transformer-based models are fine-tuned using a standardized and
conservative setup, in line with common practices for document
classification.
Given the large training set and the use of pre-trained language
models, the number of training epochs is intentionally kept small to
avoid overfitting and unnecessary computational overhead.

Models are trained for 2--4 epochs, which are sufficient to reach
convergence while preserving the generalization capabilities of the
pre-trained representations. Learning rates are selected from
commonly adopted ranges for transformer fine-tuning to ensure stable
optimization.

Batch size is fixed according to hardware constraints and is not
treated as a tunable parameter. To account for stochastic effects,
each configuration is evaluated across multiple random seeds.

Overall, hyperparameter choices for transformer models prioritize
stability and reproducibility rather than aggressive performance
optimization, in contrast with the more interpretable threshold-based
tuning of the Two-Stage classifier.
To explicitly assess the robustness of transformer fine-tuning choices,
we analyze both intra-seed variability and cross-architecture behavior
under identical training configurations.

\paragraph{Intra-seed Stability}
To evaluate the impact of stochasticity in optimization, each transformer
architecture is fine-tuned using multiple random seeds while keeping all
hyperparameters fixed.
Results show minimal variance in macro-averaged F1 across seeds,
indicating stable convergence and limited sensitivity to initialization.

\textbf{[INSERT TABLE: Intra-seed performance stability for transformer models.
Rows: model architectures; columns: different seeds, mean, standard deviation.]}

\paragraph{Cross-Architecture Hard-Case Analysis}
Beyond seed variability, we investigate whether different transformer
architectures resolve ambiguous samples in complementary ways.
To this end, we analyze prediction overlap restricted to a subset of
hard cases, defined as samples consistently misclassified by simpler
models.

The resulting overlap patterns reveal a high degree of agreement across
architectures, suggesting that residual errors are driven by intrinsic
semantic ambiguity rather than suboptimal hyperparameter choices or
model-specific artifacts.

\textbf{[INSERT FIGURE: Cross-architecture prediction overlap on hard cases
(e.g., heatmap or matrix plot across transformer models).]}

Overall, these analyses confirm that transformer fine-tuning operates in
a stable regime and that performance differences are primarily determined
by architectural inductive biases rather than aggressive hyperparameter
optimization. Consequently, transformer models are treated as robust
semantic baselines, whose complementarity is later exploited through
ensemble analysis in the Results section.
